\section{安全与安保目标}

\subsection{CIA 三要素}

\begin{importantbox}{Agentic 系统的安全目标：CIA}
\begin{itemize}
    \item \textbf{Confidentiality（机密性）}：信息仅对授权实体可访问，包括系统秘密、用户凭证、用户数据、模型本身
    \item \textbf{Integrity（完整性）}：系统和数据未被篡改，保持准确和可信
    \item \textbf{Availability（可用性）}：授权用户可靠、及时地访问数据、系统和资源
\end{itemize}
\end{importantbox}

\subsection{与传统系统的对比}

相比传统计算机系统，Agentic AI 系统有\textbf{额外的保护目标}：
\begin{itemize}
    \item \textbf{机密性}：需额外保护 API 密钥、秘密 Prompt、交互历史、专有模型参数等
    \item \textbf{完整性}：需确保模型完整性（model integrity）
    \item \textbf{可用性}：需确保模型性能和服务可用性
\end{itemize}

\begin{warningbox}{LLM 引入的攻击面增大}
LLM 的使用显著增大了攻击面：
\begin{itemize}
    \item 机密性：LLM 处理敏感信息后可能通过输出泄露
    \item 完整性：Prompt 可利用不可信数据；模型可能被投毒或包含后门
    \item 可用性：模型本身可遭受拒绝服务攻击
\end{itemize}
\end{warningbox}

